\section{\texorpdfstring{$G$}{G}-injectivity and group averaging}%
\label{sec:peps_quantum_double}

Let $G$ be a finite group. The quantum double model of $G$
\cite{Kitaev2003FaultTolerant} carries a spin with basis
$\{\ket{g}\}_{g\in G}$ on every edge of a square lattice; for
$G=\mathbb{Z}_2$ it is the toric code. The review gives two PEPS
representations \cite[Appendix~A, ``The Toric Code and quantum double
    models'']{Cirac2021Matrix}. Their symmetries are phrased through the map
$\mathcal{P}(A)=\sum A^i_{\alpha\beta\gamma\delta}\ket{i}\bra{\alpha\beta\gamma\delta}$
from the virtual to the physical system
\cite[Section~``MPS and PEPS'']{Schuch2010PEPS}.

\begin{definition}[$G$-injective and $G$-isometric tensors]%
    \label{def:peps_g_injective}
    \lean{TNLean.PEPS.IsGInjective, TNLean.PEPS.IsGIsometric,
        TNLean.PEPS.siteMap, TNLean.PEPS.siteMap_apply}
    \leanok
    Let $g\mapsto U_g$ be a representation of a finite group $G$ on the
    virtual system of a tensor $A$, and let $\Pi_U$ be the projector onto the
    $U_g$-invariant subspace. The tensor is $G$-injective if
    $\mathcal{P}(A)U_g=\mathcal{P}(A)$ for every $g$ and $\mathcal{P}(A)$ is
    injective on the invariant subspace
    \cite[Section~``Two dimensions: $G$-injective PEPS'']{Schuch2010PEPS}.
    It is $G$-isometric if in addition
    $\mathcal{P}(A)$ restricted to the invariant subspace is a multiple of an
    isometry, $\langle\mathcal{P}(A)x,\mathcal{P}(A)y\rangle
        =c\,\langle x,y\rangle$ for invariant $x,y$ and one constant $c>0$
    \cite[Section~``Isometric PEPS'']{Schuch2010PEPS}.
    When the virtual action factors as
    $U_g=U_g^{\ell}\otimes U_g^r\otimes U_g^t\otimes U_g^b$,
    the invariance of a four-leg tensor reads
    % Formula: P(A)(U_g^ell tensor U_g^r tensor U_g^t tensor U_g^b)=P(A).
    % Source: Papers/1001.3807/paper_v3.tex:1278--1288, figs3/ug-sym.
    % Ink-to-index: each of the four cardinal rays is one virtual leg;
    % its inserted matrix is the representation acting on that leg, including
    % its contragredient when dictated by the bond orientation. The northeast
    % ray is the physical index i. No physical index is suppressed.
    % Contraction: each U matrix contracts its inward port with A and exposes
    % its outward port. Both sides have (V,P)=(4,1).
    \begin{tenkzequation}
        \tenkzkernel
        \begin{tenkz}[rows={wire,wire,wire}, cols=3, bonds=none]
            \tn[at=(2,2), name=symmetricSite, label pos=225,
                ports={180:virtual, 0:virtual, 90:virtual, 270:virtual,
                        45:physical:$i$}]{A}
            \tn[at=(2,1), name=symmetryLeft, skin=ring, label pos=270,
                ports={180:virtual, 0:virtual}]{U_g^{\ell}}
            \tn[at=(2,3), name=symmetryRight, skin=ring, label pos=270,
                ports={180:virtual, 0:virtual}]{U_g^r}
            \tn[at=(1,2), name=symmetryTop, skin=ring, label pos=180,
                ports={90:virtual, 270:virtual}]{U_g^t}
            \tn[at=(3,2), name=symmetryBottom, skin=ring, label pos=0,
                ports={90:virtual, 270:virtual}]{U_g^b}
            \tnwire{symmetryLeft.0}{symmetricSite.180}
            \tnwire{symmetryRight.180}{symmetricSite.0}
            \tnwire{symmetryTop.270}{symmetricSite.90}
            \tnwire{symmetryBottom.90}{symmetricSite.270}
        \end{tenkz}
        $=$
        \begin{tenkz}[rows={wire}, cols=1, bonds=none]
            \tn[label pos=225,
                ports={180:virtual, 0:virtual, 90:virtual, 270:virtual,
                        45:physical:$i$}]{A}
        \end{tenkz}
    \end{tenkzequation}
    The four factors include the contragredient actions required by the
    orientations of the virtual bonds.
\end{definition}

\begin{lemma}[Left-inverse form of $G$-injectivity]%
    \label{thm:peps_g_injective_left_inverse}
    \lean{TNLean.PEPS.apply_averageMap_of_forall_comp_eq,
        TNLean.PEPS.isGInjective_iff_exists_leftInverse}
    \leanok
    \uses{def:peps_g_injective}
    A tensor is $G$-injective if and only if
    $\mathcal{P}(A)U_g=\mathcal{P}(A)$ for every $g$ and $\mathcal{P}(A)$ has a
    left inverse $L$ with $L\,\mathcal{P}(A)=\Pi_U$
    \cite[Section~``Two dimensions: $G$-injective PEPS'']{Schuch2010PEPS}.
\end{lemma}

\begin{proof}\leanok
    An invariant map satisfies $\mathcal{P}(A)\Pi_U=\mathcal{P}(A)$, since
    $\Pi_U$ is the average of the $U_g$. If $\mathcal{P}(A)$ is injective on
    the invariant subspace, choose a linear map $L$ into the invariant
    subspace with $L\,\mathcal{P}(A)x=x$ for invariant $x$; then
    $L\,\mathcal{P}(A)=L\,\mathcal{P}(A)\Pi_U=\Pi_U$. Conversely, if
    $L\,\mathcal{P}(A)=\Pi_U$ and $\mathcal{P}(A)x=0$ for an invariant $x$, then
    $x=\Pi_Ux=L\,\mathcal{P}(A)x=0$.
\end{proof}

\begin{lemma}[$G$-injectivity for the trivial representation]%
    \label{thm:peps_g_injective_trivial}
    \lean{TNLean.PEPS.isGInjective_trivial_iff,
        TNLean.PEPS.siteMap_injective_iff}
    \leanok
    \uses{def:peps_g_injective}
    For the trivial representation, $G$-injectivity is injectivity of
    $\mathcal{P}(A)$, that is, linear independence of the physical vectors
    $A^{\bullet}_{\alpha\beta\gamma\delta}$ indexed by the virtual
    configurations.
\end{lemma}

\begin{proof}\leanok
    Every vector is invariant under the trivial representation, and
    $\mathcal{P}(A)$ is the linear map whose columns are the vectors
    $A^{\bullet}_{\alpha\beta\gamma\delta}$.
\end{proof}

The group-theoretic input to $G$-injectivity is the averaging over the group
and an operator $\Delta$ that inverts the expansion of an operator in the
$U_g$'s \cite[Section~``MPS: the $G$-injective case'']{Schuch2010PEPS}.
Throughout, $g\mapsto U_g$ is a representation of a finite group $G$ on a
finite-dimensional complex space.

\begin{lemma}[Group averaging]%
    \label{thm:peps_group_averaging}
    \lean{Representation.averageMap_linHom_apply,
        Representation.mem_invariants_linHom_iff,
        Representation.isProj_averageMap_linHom,
        Representation.averageMap_linHom_apply_of_adjoint,
        Representation.trace_adjoint_comp_averageMap_linHom}
    \leanok
    The map $\sigma(X)=\frac{1}{|G|}\sum_g U_gXU_g^{-1}$ is a projection onto
    the operators commuting with every $U_g$. If the representation is unitary,
    then $\sigma(X)=\frac{1}{|G|}\sum_g U_gXU_g^\dagger$ and
    $\tr[Y^\dagger\sigma(X)]=\tr[\sigma(Y)^\dagger X]$, so $\sigma$ is the
    orthogonal projector onto these operators for the Hilbert--Schmidt inner
    product \cite[Section~``MPS: the $G$-injective case'']{Schuch2010PEPS}.
\end{lemma}

\begin{proof}\leanok
    The map $\sigma$ is the average of the conjugation action
    $X\mapsto U_gXU_g^{-1}$, so it is a projection onto the operators fixed
    by every conjugation, and these are the operators commuting with every
    $U_g$. For a unitary representation,
    $\tr[Y^\dagger U_gXU_g^\dagger]=\tr[(U_g^\dagger YU_g)^\dagger X]$, and
    summing over $g$, with $g$ replaced by $g^{-1}$, gives the symmetry of
    $\sigma$.
\end{proof}

\begin{definition}[Isotypic projectors and the operator $\Delta$]%
    \label{def:peps_delta_operator}
    \lean{Representation.charProjector,
        Representation.irreducibleCharacters,
        Representation.characterMultiplicity,
        Representation.deltaOperator}
    \leanok
    Let $\chi_U$ be the character of $U$, and let $\chi_1,\dots,\chi_r$ be the
    distinct characters of the irreducible subrepresentations of $U$, with
    $d_i=\chi_i(1)$ and multiplicities
    $m_i=\frac{1}{|G|}\sum_g\chi_U(g)\chi_i(g^{-1})$. The isotypic projectors
    are $P_i=\frac{d_i}{|G|}\sum_g\chi_i(g^{-1})U_g$, and
    \begin{align*}
        \Delta=\frac{1}{|G|}\sum_{i=1}^{r}\frac{d_i}{m_i}P_i
    \end{align*}
    \cite[Section~``MPS: the $G$-injective case'']{Schuch2010PEPS}.
\end{definition}

\begin{lemma}[Isotypic projectors from characters]%
    \label{thm:peps_isotypic_projectors}
    \lean{Representation.sum_character_inv_smul_eq,
        Representation.exists_isInternal_isAtom,
        Representation.charProjector_apply_of_mem,
        Representation.irreducibleCharacters_finite,
        Representation.characterMultiplicity_eq_finrank,
        Representation.sum_charProjector_irreducibleCharacters,
        Representation.trace_comp_charProjector}
    \leanok
    \uses{def:peps_delta_operator}
    For irreducible representations $D^i$ and $D^j$ of $G$,
    $\sum_g\chi_i(g^{-1})D^j(g)$ equals $\frac{|G|}{d_i}\Id$ if $D^i$
    and $D^j$ are equivalent and $0$ otherwise. Consequently $P_i$ acts as
    the identity on every irreducible subrepresentation of $U$ with
    character $\chi_i$ and as zero on every other irreducible
    subrepresentation, the number $m_i$ is the dimension of the space of
    intertwining maps from $D^i$ to $U$, and
    \begin{align*}
        \sum_{i=1}^{r}P_i & =\Id, & \tr[U_kP_i] & =m_i\chi_i(k)
    \end{align*}
    \cite[Section~``MPS: the $G$-injective case'']{Schuch2010PEPS}.
\end{lemma}

\begin{proof}\leanok
    The operator $\sum_g\chi_i(g^{-1})D^j(g)$ commutes with every $D^j(h)$,
    because characters are constant on conjugacy classes, so by Schur's lemma
    it is a multiple of the identity; its trace is
    $\sum_g\chi_i(g^{-1})\chi_j(g)$, which orthogonality of characters
    evaluates. By Maschke's theorem $U$ is a direct sum of irreducible
    subrepresentations. On each summand the projectors act as just
    described, which gives $\sum_iP_i=\Id$, and the trace of
    $U_kP_i$ is the sum of $\chi_i(k)$ over the summands with character
    $\chi_i$. Orthogonality of characters counts these summands as $m_i$.
\end{proof}

\begin{lemma}[Block form of $\Delta$]%
    \label{thm:peps_delta_block_form}
    \lean{Representation.deltaOperator_apply_of_mem}
    \leanok
    \uses{def:peps_delta_operator}
    On every irreducible subrepresentation of $U$ with character $\chi_i$,
    $\Delta$ acts as the scalar $\frac{d_i}{m_i|G|}$. Thus, in a basis in which
    $U_g\cong\bigoplus_iD^i(g)\otimes\Id_{m_i}$,
    \begin{align*}
        \Delta\cong\frac{1}{|G|}\bigoplus_i\frac{d_i}{m_i}\,
        \Id_{d_i}\otimes\Id_{m_i}
    \end{align*}
    \cite[Section~``MPS: the $G$-injective case'']{Schuch2010PEPS}.
\end{lemma}

\begin{proof}\leanok
    \uses{thm:peps_isotypic_projectors}
    Only the projector of $\chi_i$ acts nontrivially on such a
    subrepresentation, and it acts as the identity.
\end{proof}

\begin{lemma}[Expansion in the representation]%
    \label{thm:peps_delta_expansion}
    \lean{Representation.trace_comp_deltaOperator,
        Representation.sum_trace_comp_deltaOperator_smul,
        Representation.sum_trace_inv_comp_deltaOperator_comp_smul}
    \leanok
    \uses{def:peps_delta_operator}
    For every $g\in G$,
    \begin{align*}
        \sum_h\tr[U_{h^{-1}}\Delta U_g]\,U_h
        =\sum_h\tr[U_{gh^{-1}}\Delta]\,U_h=U_g .
    \end{align*}
    For a unitary representation $U_{h^{-1}}=U_h^\dagger$
    \cite[Section~``MPS: the $G$-injective case'']{Schuch2010PEPS}.
\end{lemma}

\begin{proof}\leanok
    \uses{thm:peps_isotypic_projectors}
    By the trace formula for the projectors,
    $\tr[U_k\Delta]=\frac{1}{|G|}\sum_i d_i\chi_i(k)$, so
    $\sum_k\tr[U_{k^{-1}}\Delta]U_k=\sum_iP_i=\Id$. Substituting
    $h=kg$ gives
    $\sum_h\tr[U_{gh^{-1}}\Delta]U_h=\big(\sum_k\tr[U_{k^{-1}}\Delta]U_k\big)U_g
        =U_g$, and the first expression equals the second by cyclicity of
    the trace.
\end{proof}

\begin{definition}[Semi-regular representation]%
    \label{def:peps_semiregular}
    \lean{Representation.IsSemiRegular,
        Representation.isSemiRegular_leftRegular}
    \leanok
    A representation of $G$ is semi-regular if it contains every irreducible
    representation of $G$, that is, if every irreducible representation
    admits a nonzero intertwining map into it
    \cite[Section~``MPS: the $G$-injective case'']{Schuch2010PEPS}. The
    left-regular representation is semi-regular.
\end{definition}

\begin{lemma}[Semi-regularity under representation equivalence]%
    \label{thm:peps_semiregular_equivalence}
    \lean{Representation.IsSemiRegular.of_equiv,
        Representation.Equiv.isSemiRegular_iff}
    \leanok
    \uses{def:peps_semiregular}
    Equivalent representations $\rho$ and $\sigma$ are semi-regular
    simultaneously. Thus the occurrence condition of
    \cite[Definition~4.5]{Schuch2010PEPS} is independent of coordinates.
\end{lemma}

\begin{proof}\leanok
    If $e$ intertwines $\rho$ with $\sigma$ and is invertible, every
    nonzero intertwiner $f$ from an irreducible representation into
    $\rho$ gives a nonzero intertwiner $ef$ into $\sigma$. Applying the
    same argument to $e^{-1}$ proves the converse.
\end{proof}

\begin{lemma}[Linear independence of a semi-regular representation]%
    \label{thm:peps_semiregular_trace}
    \lean{Representation.character_leftRegular,
        Representation.deltaOperator_leftRegular,
        Representation.sum_irreducibleCharacters_mul_of_isSemiRegular,
        Representation.trace_inv_comp_comp_deltaOperator_of_isSemiRegular}
    \leanok
    \uses{def:peps_delta_operator, def:peps_semiregular}
    If $U$ is semi-regular, then $\tr[U_{g^{-1}}U_h\Delta]=\delta_{g,h}$ for
    all $g,h\in G$, with $U_{g^{-1}}=U_g^\dagger$ for a unitary
    representation. For the left-regular representation
    $\Delta=\frac{1}{|G|}\Id$
    \cite[Section~``MPS: the $G$-injective case'']{Schuch2010PEPS}.
\end{lemma}

\begin{proof}\leanok
    \uses{thm:peps_delta_expansion, thm:peps_isotypic_projectors}
    The character of the left-regular representation is $|G|\delta_{k,1}$,
    so every character $\chi_i$ has multiplicity $d_i$ in it, and
    $\Delta=\frac{1}{|G|}\sum_iP_i=\frac{1}{|G|}\Id$. Every
    irreducible representation occurs in the left-regular representation,
    and, since $U$ is semi-regular, in $U$; hence $U$ and the left-regular
    representation have the same irreducible characters. Comparing the two
    expressions for $\tr[U_k\Delta]$ in the left-regular representation gives
    $\sum_id_i\chi_i(k)=|G|\delta_{k,1}$, and therefore
    $\tr[U_{g^{-1}}U_h\Delta]=\frac{1}{|G|}\sum_id_i\chi_i(g^{-1}h)
        =\delta_{g,h}$.
\end{proof}

\begin{lemma}[Semi-regularity and the group algebra]
    \label{thm:peps_semiregular_group_algebra}
    \lean{Representation.linearIndependent_iff_injective_asAlgebraHom,
        Representation.isSemiRegular_of_linearIndependent,
        Representation.isSemiRegular_iff_linearIndependent,
        Representation.isSemiRegular_iff_injective_asAlgebraHom}
    \leanok
    \uses{def:peps_semiregular}
    Let $G$ be a finite group and let $\rho$ be a finite-dimensional complex
    representation. The following assertions are equivalent:
    $\rho$ contains every irreducible representation; the operators
    $(\rho(g))_{g\in G}$ are linearly independent; and the induced algebra
    homomorphism $\mathbb C[G]\to\operatorname{End}(V)$ is injective.
    This is an algebraic characterization of the occurrence condition in
    \cite[Definition~4.5 and Lemma~4.6]{Schuch2010PEPS}.
    More generally, for a representation of a monoid over a commutative
    semiring, independence of its operators is equivalent to injectivity
    of its monoid-algebra homomorphism.
\end{lemma}
\begin{proof}\leanok
    \uses{thm:peps_semiregular_trace}
    The trace-dual identity gives independence for a semi-regular
    representation. Conversely, if an irreducible character $\chi$ is
    absent, its character projector vanishes:
    \begin{align}
        \frac{\chi(1)}{|G|}\sum_{g\in G}\chi(g^{-1})\rho(g)=0.
        \label{eq:peps_absent_character}
    \end{align}
    Indeed, the projector vanishes on every irreducible summand of $V$.
    Independence in (\ref{eq:peps_absent_character}) would imply
    $\chi(1)=0$, contrary to the positive dimension of the irreducible
    representation. The standard basis of the group algebra identifies
    the last assertion with independence of the group operators.
\end{proof}

\begin{lemma}[Reversing a semi-regular bond]
    \label{thm:peps_semiregular_dual}
    \lean{Representation.IsSemiRegular.dual,
        Representation.isSemiRegular_dual_iff}
    \leanok
    \uses{def:peps_semiregular}
    A finite-dimensional complex representation $\rho$ of a finite group
    is semi-regular if and only if its contragredient
    $\rho^*(g)=\rho(g^{-1})^{\mathsf T}$ is semi-regular.
    This justifies reversing the bond orientation in the contraction
    convention of \cite[Definition~5.1 and Lemma~5.2]{Schuch2010PEPS}.
\end{lemma}
\begin{proof}\leanok
    \uses{thm:peps_semiregular_group_algebra}
    Inversion permutes the group indices, and transposition is injective
    on endomorphisms of a finite-dimensional vector space. It therefore
    preserves independence of the group operators. Conversely,
    independence of their transposes implies independence of the
    original operators.
\end{proof}

\begin{lemma}[Nonzero tensor factors of a semi-regular representation]
    \label{thm:peps_semiregular_nonzero_tensor}
    \lean{Representation.IsSemiRegular.tprod_of_nontrivial,
        Representation.IsSemiRegular.tprod_left_of_nontrivial}
    \leanok
    \uses{def:peps_semiregular}
    Let $G$ be a finite group, let $\rho$ be a finite-dimensional
    semi-regular complex representation, and let $\sigma$ be any
    nonzero finite-dimensional complex representation. Then both
    $\rho\otimes\sigma$ and $\sigma\otimes\rho$ are semi-regular.
    This strengthens the tensor-product assertion in
    \cite[Lemma~5.2]{Schuch2010PEPS}: the second factor need not contain
    the trivial representation.
\end{lemma}
\begin{proof}\leanok
    \uses{thm:peps_semiregular_group_algebra, thm:peps_semiregular_trace,
        thm:peps_semiregular_equivalence}
    A dependence among the tensor-product operators is a dependence
    among the tensors $\rho(g)\otimes\sigma(g)$, by the natural
    identification of their endomorphism spaces. Applying the
    trace-dual functional at $h$ to the first factor gives
    $c_h\sigma(h)=0$. Since $\sigma(h)$ is invertible on a nonzero
    vector space, $c_h=0$. The group-algebra criterion proves
    semi-regularity. Interchanging the tensor factors gives the
    other order.
\end{proof}

\begin{lemma}[Irreducibility of the contragredient]
    \label{thm:peps_irreducible_dual}
    \lean{Representation.isIrreducible_dual}
    \leanok
    The contragredient of a finite-dimensional irreducible complex
    representation of a group is irreducible. No unitarity or finite-group
    hypothesis is needed. This is an algebraic ingredient for the
    conjugate charge in \cite{Schuch2010PEPS}, lines 2514--2518.
\end{lemma}
\begin{proof}\leanok
    The annihilator of an invariant dual subspace is invariant in the
    original representation, since the group action on the dual uses
    the inverse element. Irreducibility makes this annihilator zero or
    the whole space. Finite-dimensional annihilator duality then makes
    the original dual subspace the whole space or zero, respectively.
\end{proof}


\begin{lemma}[Tensor products of semi-regular representations]%
    \label{thm:peps_semiregular_tensor_product}
    \lean{Representation.isIrreducible_of_finrank_eq_one,
        Representation.IsSemiRegular.exists_mem_invariants_ne_zero,
        Representation.IsSemiRegular.tprod_of_mem_invariants,
        Representation.IsSemiRegular.tprod}
    \leanok
    \uses{def:peps_semiregular}
    If $U$ and $V$ are semi-regular representations of $G$, then so is
    $U\otimes V$
    \cite[Section~``Two dimensions: $G$-injective PEPS'']{Schuch2010PEPS}.
\end{lemma}

\begin{proof}\leanok
    A one-dimensional representation is irreducible, so the semi-regular
    representation $V$ contains the trivial representation and has a nonzero
    invariant vector $v$. If $D$ is irreducible and $f$ is a nonzero
    intertwining map from $D$ to $U$, then $x\mapsto f(x)\otimes v$ is a
    nonzero intertwining map from $D$ to $U\otimes V$.
\end{proof}
